\section{Ein Satz im Ziel: DPrcxq66 erfüllt im Gesamtmodell alle Grenzen}\label{sec:ziel}

\textbf{Lesehilfe.} „Im Ziel“ heißt hier: Das Gesamtmodell (alle Spulen, Mittel über die Spulen) hält die angenommenen Grenzen ein.
Die Grenzen sind Annahmen und keine gemessenen Werte: Stahl-Vergleichsspannung (p99, 99\,\% des Materials liegen darunter) höchstens 500\,MPa,
Dehnung des Wickelpakets höchstens 0,3\,\%, Neutronen-Fluenz höchstens $3\cdot10^{18}$\,n/cm$^2$, Abschirmung mindestens 1,0\,m.
Die Hülle ist der freie Abstand zwischen Plasmarand und Gehäuse-Innenseite der Spule.

Der Nutzer wollte, dass wir einen Satz „ins Ziel bringen, auch um zu verstehen“, und ihn dann als Absprungbasis nehmen.

\subsection*{Weg dorthin}
\begin{enumerate}\itemsep1pt
\item Verbinder nach Lücken setzen: größte Strecke ohne Verbinder von 17,4\,m auf 4\,m, Auslenkung von 241 auf 42\,mm (DjLe3eL6).
\item Platz für die Abschirmung: Der Satz mit dem besten Plasma hat zu wenig Hülle (0,47\,m). Anlage vergrößern (DYU $\times1{,}45$) schafft Platz, aber die Spannung
  bleibt nicht gleich (417 auf 627\,MPa). Ein Satz mit mehr Abstand (DPrcxq66) war der bessere Weg.
\item Dickere Wände bei diesem Satz.
\end{enumerate}

\subsection*{Ergebnis (Gesamtmodell)}
\begin{tabular}{lrrrrl}
\hline
Variante & Stahl p99 & Dehnung & Hülle & Fluenz & alle Grenzen \\ \hline
DPrcxq66, Wände $\times1{,}3$ & 453\,MPa & 0,27\,\% & 1,32\,m & $1{,}1\cdot10^{18}$ & ja \\
DPrcxq66, Wände $\times1{,}6$ & 384\,MPa & 0,26\,\% & 1,56\,m & $1{,}1\cdot10^{17}$ & ja \\
DGHe48vB, Wände $\times1{,}3$ & 483\,MPa & 0,34\,\% & 1,35\,m & $7{,}9\cdot10^{17}$ & nein (Dehnung) \\
DYU $\times1{,}45$, größeres Lager & 500\,MPa & 0,30\,\% & 1,27\,m & $9{,}7\cdot10^{17}$ & knapp nein \\ \hline
\end{tabular}

\smallskip
\noindent DPrcxq66 hat zwei Feldperioden, QI-Güte 0,00079 (rund achtmal besser als DP9 mit 0,0068), Kraft 792\,MN je Spule.

\subsection*{Erwartung vorher und Vergleich}
Erwartet waren 460 bis 520\,MPa ($\times1{,}3$) und 400 bis 450\,MPa ($\times1{,}6$); gefunden 453 und 384\,MPa, also etwas besser. Die Dehnungsbereiche trafen,
nur DGHe48vB lag mit 0,34\,\% knapp über dem erwarteten Bereich 0,30 bis 0,33\,\%.

\subsection*{Grenzen der Aussage}
\begin{itemize}\itemsep1pt
\item Die Grenzen sind Annahmen. Das Submodell der heißesten Spule liefert 602 bis 605\,MPa, also über 500\,MPa. Welcher Maßstab gelten soll, ist offen.
\item Hülle 1,32\,m liegt knapp über dem belegten HELIAS-Wert von 1,3\,m, Kachel- und Toleranzwerte fehlen.
\item Abkühlen und ein feineres Netz sind für diesen Satz in Arbeit; erst danach gilt er als Absprungbasis.
\item QI ist ein Gütemaß der Feldform, Einschluss und Stabilität sind nicht gerechnet.
\end{itemize}
